\section{Actual initial identifiers on sector runs}
\label{sec:area_law_actual_sector_runs}

Fix an arbitrary origin $o\in\mathbb R^2$, a finite endpoint set
$Z\subset\mathbb Z^2$, an integer width $C\ge2$ and an initial layer
$k_0\ge50{,}000{,}000$. Choose arbitrary valid pitch residues at every
layer. Let $\mathcal I$ be the actual initial identifiers, let $X_i^0$
be their closed birth regions, put $U_i=\operatorname{int}(X_i^0)$,
and write $\chi(i)\in\{0,1\}$ for their initial colors.
Let $k\ge k_0$, let $z\in F_{k,\ell_k}$ be an actual fine-cell index,
and let $v$ be any of the nine marks of its defining cell.
Put $\ell=\ell_k-5$ and $r=2^\ell/2=t_k/64$.

Use the all-midpoint fans centered at $v$, with origins
$o_{\rm S}=v-r(1,1)$ and $o_{\rm L}=v-2r(1,1)$, exponents $\ell$
and $\ell+1$, and cell index $(0,0)$. Let $\mathcal J$ index their
eight corresponding elementary bases, and write $P_s^{\rm S}$ and
$P_s^{\rm L}$ for the closed triangles. Let
$\sigma:\mathcal J\to\mathcal I$ be the unique actual assignment of
Theorem~\ref{thm:area_law_actual_sector_assignment}.
For the smaller fan, write $b_s^-$ and $b_s^+$ for the initial and
final endpoints of the oriented base $s$. Two bases are adjacent when
$b_s^+=b_t^-$ or $b_t^+=b_s^-$.

Give base $s$ the color $\chi(\sigma(s))$. Form the existing fan graph
$G$, joining distinct adjacent bases exactly when their colors agree.
Let $\mathcal R$ be its connected components, and write $[s]$ for the
component of $s$. For $R\in\mathcal R$, its closed run region is
$Y_R=\bigcup_{s\in R}P_s^{\rm S}$.
These local runs occur in \cite[Section~11]{OpenAI2026AreaLaw}.

% Source: 10-geometry.tex, geometry:initial-stars, lines 333--370,
% especially 361--370, and prop:two-families, lines 313--323;
% revision adc7f1241b42e322a6451854ab7e4b4c146bf78a.
% These entries do not assert active-ray enumeration or the full star theorem.

\begin{theorem}[Initial identifiers on cyclic runs]
    \label{thm:area_law_initial_sector_run_assignment}
    \lean{TNLean.PEPS.AreaLaw.Geometry.exists_unique_initialRegion_sector_run_assignment}
    \leanok
    \uses{thm:area_law_actual_sector_assignment,
        def:area_law_initial_region_indices,
        def:area_law_initial_birth_regions,
        def:area_law_initial_region_colors,def:area_law_fan_runs,
        def:area_law_cell_fan}
    There is a unique function $\lambda:\mathcal R\to\mathcal I$
    satisfying $\lambda([s])=\sigma(s)$ for every $s\in\mathcal J$
    and the following identity for every $i\in\mathcal I$:
    \begin{align}
        X_i^0\cap\overline B_\infty(v,r)
         & =\bigcup_{\substack{R\in\mathcal R \\\lambda(R)=i}}Y_R.
    \end{align}
    Distinct components carrying one identifier remain distinct terms
    of the union. If no sector has identifier $i$, the union is empty.
\end{theorem}
\begin{proof}\leanok
    \uses{thm:area_law_initial_sector_colors,
        thm:area_law_actual_sector_assignment,def:area_law_fan_runs}
    An edge of $G$ joins adjacent sectors with equal colors.
    The adjacent-sector color theorem therefore gives equality of
    their assigned identifiers. Equality is preserved along paths,
    so $\sigma$ is constant on every connected component. It thus
    determines $\lambda([s])=\sigma(s)$. Every component contains a
    sector, which also makes this function unique.

    The actual closed restriction is
    $X_i^0\cap\overline B_\infty(v,r)
        =\bigcup_{\sigma(s)=i}P_s^{\rm S}$.
    Every sector belongs to its component, and all sectors in a
    component labeled $i$ have assigned identifier $i$.
    Grouping the terms by their connected components gives the
    asserted identity. The grouping retains every separated component
    carrying the same identifier. When there are no such components,
    both unions are empty.
\end{proof}

\begin{theorem}[The case with no change of identifier]
    \label{thm:area_law_initial_sector_no_active}
    \lean{TNLean.PEPS.AreaLaw.Geometry.exists_unique_initialRegion_of_no_active_sector}
    \leanok
    \uses{thm:area_law_actual_sector_assignment,
        def:area_law_initial_region_indices,
        def:area_law_initial_open_regions,def:area_law_cell_fan}
    Suppose that $\sigma(s)=\sigma(t)$ for every adjacent pair
    $s,t\in\mathcal J$. Then there is a unique identifier
    $i\in\mathcal I$ such that
    $\overline B_\infty(v,r)\subseteq U_i$.
\end{theorem}
\begin{proof}\leanok
    \uses{thm:area_law_fan_run_monochromatic,
        thm:area_law_initial_sector_run_assignment,
        thm:area_law_actual_sector_assignment,
        thm:area_law_cell_fan_regularity,
        thm:area_law_initial_birth_regularity,
        thm:area_law_cell_fan_cover,
        thm:area_law_centered_dyadic_cell_ball,
        thm:area_law_initial_open_disjoint,def:area_law_fan_runs}
    The constant-colored fan graph is connected, including across
    the last and first bases. The assumption makes it a subgraph of
    $G$: neighboring equal identifiers have equal initial colors.
    Hence $G$ has one connected component, and descent shows that
    $\sigma$ is constant at some actual identifier $i_0$.

    For each working triangle,
    $\operatorname{int}(P_s^{\rm L})\subseteq U_{i_0}$.
    Taking closures and using triangle and birth-region regularity
    gives $P_s^{\rm L}\subseteq X_{i_0}^0$. The actual fan cover and
    the centered-cell identity therefore yield
    \begin{align}
        \overline B_\infty(v,2r)
         & =\bigcup_{s\in\mathcal J}P_s^{\rm L}
        \subseteq X_{i_0}^0.
    \end{align}
    Taking interiors gives $B_\infty(v,2r)\subseteq U_{i_0}$.
    Since $r>0$, the smaller closed square lies in this larger open
    square, and thus $\overline B_\infty(v,r)\subseteq U_{i_0}$.
    Finally $v$ belongs to the smaller square. Pairwise disjointness
    of the actual open regions prevents another identifier from
    containing that square, proving uniqueness.
\end{proof}
